% content.tex — the body of paper_04_language_of_the_gift, as a book chapter.
% GENERATED by split_content.py. The standalone paper still builds from
% paper_04_language_of_the_gift.tex; both read this file, so they cannot drift apart.

\begin{abstract}
\noindent
The gift is the native language of love: an expressive act whose meaning is
constituted by the relation it enacts. Economics possesses a rich literature
on gift-giving---altruism in the family, gift exchange in labor markets,
gifts as costly signals, the deadweight loss of in-kind transfer, and the
crowding-out of intrinsic motivation---yet each of these theories, I argue,
shares a common move: the gift is made tractable by \emph{interpreting} it
into a transfer language whose semantics is a ledger. This paper makes that
move itself the object of analysis. A semiotic and psychoanalytic layer
first explains the gift's surplus over the thing: following Saussure and
Baudrillard, the gift is a signifier whose value is positional---generated
by its place in a differential system that is the relation's own history---
and following Lacan, what it signifies is desire rather than need, so that
its functional excess (the economist's ``deadweight loss'') is the mechanism
of its signifying work, not waste. Using elementary tools from formal
language theory, I then model the expressive practice of intimate giving as a
context-sensitive language whose semantics depends on relational history,
and the economic-institutional reading as a non-injective, context-forgetting
interpretation map into a finite-state ledger language. \emph{Alienation} is
then given a precise structural characterization: it is the contraction of
generative practice onto the preimage of the ledger language---agents come
to produce only what the interpretation can see. Two further harms follow:
\emph{epistemic injustice} \parencite{fricker2007}, since contributions outside
the map's range are illegible as contributions; and \emph{distributive
misdescription}, since ledger balance is the wrong functional for value that
circulates generatively. Building on Eglash's generative justice and a
relational ontology of the gift, I propose recirculation depth, rather than
bilateral balance, as the appropriate justice functional, and draw
implications for the evidence-informed use of intimate-relationship data.
\end{abstract}
\clearpage

% 致辞 (Dedication) — its own page
\thispagestyle{empty}
\null\vfill
\begin{center}
\itshape
To the woman I love most deeply,\\
my ``girl of the forest and the mountains'':\\[4pt]
to have met her is the finest gift ever given me,\\
and the whole of my life is the gift I return.
\end{center}
\vfill\null
\clearpage

\section{Introduction}
\label{p04:sec:intro}

Within intimate relationships, gifts are utterances before they are
transfers. A gift says something---attention, recognition, commitment---whose
content is not carried by the object but by the act-in-relation. Yet whenever
giving must be made legible to institutions, or simply auditable to the
partners' own retrospective accounting, it is read through a different
system: a language of transfers, valuations, and balances. The thesis of this
paper is that the resulting harm, which I call the \emph{alienation of the
gift}, is best understood not as a vague cultural lament but as a structural
property of the \emph{interpretation} between two formal systems of unequal
expressive power---and that this analysis unifies three literatures that
rarely meet: the economics of gift-giving, the epistemology of injustice, and
the theory of generative justice.

This paper is the fourth in a series on the philosophy of intimacy and the
theory of justice. Paper III \parencite{huang2026vow} analyzed the unilateral
\emph{vow} as the institution of a rights-order and located its injustice in
the monopolization of second-order constituent power; the present paper turns
from the legislation of a shared life to its \emph{circulation}---from the
law of love to the language of love---and finds a structurally parallel
diagnosis: just as a vow can institute domination while remaining sincere,
an accounting of gifts can institute alienation while remaining accurate.

The argument proceeds in four steps.

\begin{enumerate}[label=(\roman*)]
  \item Section~\ref{p04:sec:econ} reviews the principal economic theories of the
  gift and identifies their shared formal move: tractability is purchased by
  interpreting expressive acts into a transfer calculus.
  \item Section~\ref{p04:sec:semiotic} opens the black box of expressive content.
  Drawing on structural semiotics and psychoanalysis, it locates the gift's
  surplus over the thing in \emph{sign-value}---value generated by the gift's
  position in a differential system that is the relation's own history---and
  identifies what the gift signifies: desire rather than need, so that the
  gift's functional excess (the economist's ``deadweight loss'') is
  constitutive rather than wasteful.
  \item Section~\ref{p04:sec:formal} formalizes the move with formal-language
  tools: an expressive language $\LE$ with context-dependent semantics, a
  ledger language $\LX$ with balance semantics, and a context-forgetting
  interpretation map $h:\LE\to\LX$. Alienation is characterized as the
  contraction of practice onto $h^{-1}(\LX)$ under legibility pressure.
  \item Section~\ref{p04:sec:epistemic} shows that the kernel of $h$---the
  expressive distinctions the map collapses---is exactly where
  \textcite{fricker2007}'s testimonial and hermeneutical injustices arise in the
  intimate sphere.
  \item Sections~\ref{p04:sec:generative}--\ref{p04:sec:data} replace the balance
  functional with a generative one (recirculation within the value-creating
  relation, after \cite{eglash2016}) and draw design implications for
  relational data systems.
\end{enumerate}

% ------------------------------------------------------------
\section{Economic theories of the gift}
\label{p04:sec:econ}

The economics of the gift is not a blind spot but a developed field. The
point of this section is not that economists ignore gifts; it is that each
theory renders the gift tractable through a characteristic \emph{reduction},
and that the reductions, though different in content, are identical in form.

\subsection{Altruism and the family: the gift as interdependent utility}
\textcite{becker1981} models intra-family transfers through interdependent
utility: the giver's utility takes the recipient's consumption or utility as
an argument, so a gift is an allocation chosen by a (possibly altruistic)
head maximizing a household objective. The Rotten Kid theorem and its
descendants show how transfers discipline behavior inside the family. The
reduction: the gift's content is exhausted by its effect on the
\emph{allocation}; the act of giving has no semantic remainder beyond the
resources moved and the preferences revealed.

\subsection{Gift exchange: the gift as implicit contract}
\textcite{akerlof1982} famously reads above-market wages and above-minimum
effort as a \emph{partial gift exchange} between firm and worker: each side
gives more than required, sustained by norms of reciprocity. Subsequent
experimental work \parencite{fehr1993} confirmed reciprocity as a robust
behavioral force. The reduction: the gift is one leg of a deferred bilateral
exchange; its meaning is its position in a reciprocity schedule. What makes
gift exchange ``work'' is precisely that the gift creates an obligation---
which is to say, an entry in an implicit ledger.

\subsection{Signaling: the gift as costly information transmission}
\textcite{camerer1988} models gifts as economic signals of the giver's type or
of commitment to the relationship: inefficient, hard-to-resell gifts are
\emph{good} signals exactly because they are costly and relationship-specific.
\textcite{prendergast2001} develop the non-monetary character of gifts as an
information problem: cash would be efficient but uninformative about how well
the giver knows the recipient. The reduction: expressive content survives,
but only as \emph{information about a hidden parameter}---the gift means
something solely insofar as it separates types in equilibrium.

\subsection{Efficiency: the deadweight loss of giving}
\textcite{waldfogel1993} estimates that in-kind gifts destroy between a tenth
and a third of their purchase value, since recipients value the goods below
cost---the ``deadweight loss of Christmas.'' The reduction is here at its
purest: the gift is evaluated as a resource transfer and found wanting, the
expressive remainder appearing, if at all, as an unmodeled ``sentimental
value'' residual. Section~\ref{p04:sec:lacan} will return to this finding and
invert it.

\subsection{Crowding out: the gift destroyed by the price}
A fourth strand shows what happens when the two registers collide.
\textcite{titmuss1970} argued that paying for blood donation degrades both the
supply and the meaning of the donated blood; \textcite{gneezy2000} showed that
introducing a fine for late pickup at daycare \emph{increased} lateness,
because the fine re-coded a moral relation as a priced service; and
\textcite{benabou2006} provide the canonical model in which extrinsic incentives
crowd out intrinsic and reputational motivation. This literature is the
closest economics comes to the present thesis: it demonstrates empirically
that \emph{re-description in the transfer language changes the practice
described}. But it lacks a general account of why; Section~\ref{p04:sec:formal}
supplies one.

\subsection{Relational earmarking: the boundary case}
Finally, economic sociology---especially \textcite{zelizer2005}---documents how
intimates actively \emph{mark} money and goods to keep relational meanings
distinct: allowances, treats, support payments are differentiated precisely
to prevent the collapse of relation into transaction. Zelizer's ``relational
work'' can be read, in the terms of this paper, as agents' defensive
maintenance of an expressive grammar against an encroaching interpretation
map.

\subsection{The shared move}
\label{p04:sec:sharedmove}
Across these theories, the gift becomes tractable when its content is
re-expressed in a language of transfers: an allocation
(\S2.1), a ledger entry in a reciprocity schedule (\S2.2), a signal
about a parameter (\S2.3), a resource flow with a welfare loss (\S2.4),
a priced incentive (\S2.5). Each reduction is legitimate for its purpose, and
several are empirically powerful. The claim defended below is not that the
reductions are false but that they are \emph{interpretations}: structure-
forgetting maps from a richer system into a poorer one---and that when the
interpretation becomes the operative description inside the relationship
itself, the forgetting becomes a harm.

% ------------------------------------------------------------
\section{The semiotics of the gift: the surplus over the thing}
\label{p04:sec:semiotic}

Section~\ref{p04:sec:sharedmove} concluded that the economic theories are
interpretations---structure-forgetting maps from a richer system into a
poorer one. Before formalizing the map (Section~\ref{p04:sec:formal}), we must
say what, positively, the richer system contains: \emph{where} the gift's
value-beyond-the-thing comes from, and \emph{what} the gift signifies. Two
traditions answer, and their answers slot precisely into the gap the
economic theories leave open.

\subsection{Sign-value: the gift as a position in a system of differences}
\label{p04:sec:signvalue}
Saussure's two founding theses are the arbitrariness of the sign---the link
between signifier and signified is fixed by the system, not by nature---and
the differential theory of value: the value of a term is purely positional,
constituted by its relations of difference to the other terms of the system
\parencite{saussure1916}. Transposed to the gift: a gift signifies not by the
intrinsic properties of the object but by its position in a relational
code---\emph{these} flowers rather than those, flowers rather than jewelry, a
gift rather than none, on this day rather than another, from this giver
after that history. \textcite{baudrillard1981} drew the economic consequence:
beyond use-value and exchange-value there is \emph{sign-value}, the worth an
object bears as a marker of position in a code; and distinct from all three
there is \emph{symbolic exchange}, the register of the gift proper, in which
the object is not autonomous but inseparable from the concrete relation of
giver and receiver and from the moment of its giving. The ``more than the
thing'' that the gift carries is therefore neither sentimental residue nor
measurement error. It is value of a different kind, generated by the gift's
place in a differential system---and the system in question is nothing other
than the history of the relation itself.

\subsection{The signified is desire: the psychoanalytic deepening}
\label{p04:sec:lacan}
If the gift is a signifier, what does it signify? Lacan's triad of need,
demand, and desire supplies the answer \parencite{lacan2006}. A \emph{need} can
be satisfied by an object. A \emph{demand}, because it is addressed to
another, is always also a demand for love---for the other's desire---over and
above whatever object it asks for. \emph{Desire} is the remainder when need
is subtracted from demand, and no object satisfies it. The gift addresses
itself to this remainder: what it gives is not, finally, the object but the
giving---the attestation that the giver's desire is engaged with this
particular other. Hence Lacan's formula that love is giving \emph{what one
does not have}: the true gift transmits something the giver does not possess
as a good among goods---their lack, their desire---and no inventory of
transferred objects can contain it.

Two consequences follow. The first is the \emph{constitutive excess}. If the
gift's work is to signify desire rather than to satisfy need, then a degree
of functional uselessness is not a defect of the gift but its signifying
mechanism: a gift wholly absorbed by need---cash, the exact item on the
list---risks signifying nothing beyond the need it meets. This inverts the
efficiency finding of \S2.4. Waldfogel's deadweight loss, the measured gap
between the cost of the gift and the recipient's use-valuation of it, is, on
the semiotic reading, in part the \emph{footprint of sign-value}: the
portion of the act that was never addressed to use in the first place, and
that was doing the signifying. To keep the claim honest: not all measured
loss is sign-value---mismatch and giver ignorance are real, and the
signaling literature (\S2.3) captures the informational part---but a nonzero
excess is constitutive, so that a gift regime engineered to zero loss would
be one in which the gift had ceased to signify.

The second is the \emph{metonymic presence of the giver}. Mauss's informants
said the gift carries the \emph{hau}, the spirit of the giver
\parencite{mauss1925}. In semiotic terms the gift is a metonym: contiguous with
the giver, a detached part standing for the person. This is why the
destruction or careless regifting of a gift wounds in a way that no
equivalent-value substitution explains, and why the gift-object resists
fungibility: its \emph{identity}, not its value, does the work.

\subsection{What a ledger can record: referent without signifier}
\label{p04:sec:referent}
A ledger records referents: which object moved, between whom, at what
imputed value. The signifier's position in the differential system---
everything established in \S\ref{p04:sec:signvalue}--\ref{p04:sec:lacan}---is not a
property of the object and is not recordable record-by-record.
Baudrillard's thesis that political economy is \emph{founded} on the
reduction of symbolic exchange to economic exchange
\parencite{baudrillard1981} is, in the terms of this paper, the historical-scale
version of the claim the next section defends formally: the accounting
interpretation flattens the signifier into its referent and sign-value into
exchange-value. What the formalization adds is the demonstration that the
flattening is \emph{architectural}---located in the per-symbol form of the
interpretation map---and therefore cannot be repaired by more generous
imputation.

% ------------------------------------------------------------
\section{A formal-linguistic account of alienation}
\label{p04:sec:formal}

This section makes the ``shared move'' of Section~\ref{p04:sec:sharedmove}
precise using elementary formal language theory. The formalism is deliberately
minimal; its purpose is to locate \emph{where} in the interpretive pipeline
the alienation occurs, not to claim that intimate life is literally a formal
language. (Remark~\ref{p04:rem:scope} states the scope limits explicitly.)

\subsection{The expressive language}
\begin{definition}[Expressive practice]
\label{p04:def:expressive}
Let $\SE$ be an alphabet of \emph{expressive act types} (giving an object,
performing care, withholding, anticipating a need, marking an occasion,
\ldots). An intimate relationship determines a set of well-formed
\emph{histories} $\LE \subseteq \SE^{*}$ together with a meaning function
\[
\mu : \LE \longrightarrow M,
\]
where $M$ is a space of relational meanings, and $\mu$ is
\emph{history-dependent}: in general the meaning of a history $wa$ (the
history $w$ extended by act $a$) is not determined by the meaning of $a$
taken alone, i.e.\ $\mu$ does not factor through a per-symbol valuation
$\nu:\SE \to M$ with $\mu(a_1\cdots a_n) = \bigoplus_i \nu(a_i)$ for any
associative $\oplus$.
\end{definition}

The non-factorization clause is the formal content of the everyday
observation that the \emph{same} act---a bouquet, a silence, a repayment---
means different things depending on what came before, on what was foregone,
on anniversaries and apologies and inside references. Meaning in $\LE$
exhibits \emph{long-distance dependencies}: an act may take its sense from an
event arbitrarily far back in the history. In grammar-theoretic terms, any
generative description of $(\LE,\mu)$ adequate to such dependencies must
exceed the finite-state, and in general the context-free, level: the
semantics is context-sensitive in the strict sense that the contribution of a
symbol depends on its (unboundedly distant) context.\footnote{The point is
the semantic analogue of the classical expressiveness arguments from
cross-serial dependencies in natural language \parencite{shieber1985}. Nothing
below depends on the exact rung of the hierarchy, only on the strict
inclusion: ledger semantics is per-symbol compositional in a way relational
semantics is not.} Section~\ref{p04:sec:semiotic} supplied this clause's ground:
the value of the gift-signifier is positional (\S\ref{p04:sec:signvalue}), and
position in a system is precisely what no per-symbol valuation $\nu$ can
carry.

\subsection{The ledger language}
\begin{definition}[Ledger]
\label{p04:def:ledger}
Let $\SX = \{(i \to j, v) : i,j \in N,\, v \in \mathbb{R}_{\ge 0}\}$ be an
alphabet of \emph{transfer records} over a set of parties $N$, and let
$\LX = \SX^{*}$ with the \emph{balance semantics}
\[
\beta : \LX \to \mathbb{R}^{N}, \qquad
\beta(x)_k \;=\; \sum_{(i\to k,\,v)\in x} v \;-\; \sum_{(k\to j,\,v)\in x} v .
\]
\end{definition}

Two properties of $\beta$ matter. It is \emph{per-symbol compositional}
(each record contributes independently; order is irrelevant up to the
running balance), hence recognizable by finite means---this is what makes
ledgers auditable. And it is \emph{total}: every string of records has a
balance. The ledger never says ``meaningless''; it always returns a number.

\subsection{The interpretation map and its kernel}
\begin{definition}[Accounting interpretation]
\label{p04:def:interp}
An \emph{accounting interpretation} is a monoid homomorphism
$h:\SE^{*}\to\SX^{*}$ induced by a per-symbol assignment
$h:\SE\to\SX\cup\{\varepsilon\}$: each expressive act type is mapped to a
transfer record (its ``imputed value'') or erased.
\end{definition}

Every theory in Section~\ref{p04:sec:econ} implements some $h$: Becker's
allocations, Akerlof's implicit reciprocity entries, Camerer's signal costs,
Waldfogel's purchase prices. Three structural facts about any such $h$ do the
philosophical work:

\begin{enumerate}[label=(\alph*)]
\item \textbf{Context-forgetting.} Because $h$ is per-symbol, it cannot be
sensitive to history: $h(wa)=h(w)h(a)$. Whatever part of $\mu$ depends on
long-distance context is invisible to $\beta\circ h$ \emph{by construction},
not by accident of calibration. No re-weighting of imputed values repairs
this, since the failure is architectural (Definition~\ref{p04:def:expressive}'s
non-factorization clause).
\item \textbf{Non-injectivity with structured kernel.} Distinct histories
$w \neq w'$ with $h(w)=h(w')$ abound, and the collapsed distinctions are
precisely the expressive ones: an act of care and a strategic appeasement of
equal imputed value are identified. Write
$\ker h = \{(w,w') : h(w)=h(w')\}$; the harm analyzed below lives in
$\ker h$.
\item \textbf{Erasure.} Acts assigned $\varepsilon$---typically sustained,
unpriceable, anticipatory care---do not merely lose nuance; they vanish from
the record altogether. Erasure is the limiting case of (b) and the formal
hook for Section~\ref{p04:sec:epistemic}.
\end{enumerate}

\subsection{Alienation as contraction onto the preimage}
\label{p04:sec:contraction}
So far, $h$ is only a description, and a lossy description harms no one. The
harm begins when the description acquires \emph{normative force}: when
institutions (courts, tax authorities, platforms) and eventually the partners
themselves audit the relationship through $\beta \circ h$. Call this
\emph{legibility pressure}: histories are rewarded, defended, or even merely
\emph{rememberable} to the extent that their images under $h$ are
well-formed and balanced.

\begin{definition}[Alienation]
\label{p04:def:alienation}
Let $G$ be the generative process by which partners produce expressive
histories (their practice). \emph{Alienation} is the dynamic contraction of
$G$ under legibility pressure: the practiced language
$L(G) \subseteq \LE$ converges toward a sublanguage on which $\mu$ factors
through $h$---i.e.\ the partners come to produce only meanings that survive
the interpretation, so that, in the limit, $(L(G),\mu)$ is isomorphic to a
fragment of $(\LX,\beta)$.
\end{definition}

\begin{proposition}[Informal]
Under legibility pressure, the fixed point of the practice dynamics is a
ledger: expressive capacity strictly greater than finite-state is not merely
unrecorded but \emph{unexercised}.
\end{proposition}

This is the precise sense in which ``love's language is re-coded as
transaction.'' The crowding-out findings of \S2.5 are the empirical
signature of the contraction: \citeauthor{gneezy2000}'s parents did not
misunderstand the fine; their practice re-formed around the language the fine
made operative. Alienation, so characterized, is not commodification (an
object entering market circulation) but a \emph{grammar replacement}: the
generative system of the relation is substituted by the image system of its
own accounting. And Zelizer's relational earmarking (\S2.6) is its
countermeasure: the deliberate maintenance of distinctions that $h$ would
collapse---in our terms, agents acting to keep $\ker h$ semantically alive.

\begin{remark}[Scope]
\label{p04:rem:scope}
The formalism claims (i) that ledger semantics is per-symbol compositional
while relational semantics is not, and (ii) that this gap, combined with
legibility pressure, predicts contraction of practice. It does not claim that
$\SE$, $\mu$, or $G$ are empirically identifiable objects, nor that a unique
$h$ exists. The model is an explanatory skeleton---a statement of \emph{which
structural features} produce the alienation---and should be held to that
standard, no more.
\end{remark}

% ------------------------------------------------------------
\section{Epistemic injustice: the kernel made flesh}
\label{p04:sec:epistemic}

\textcite{fricker2007} distinguishes \emph{testimonial} injustice (a deflated
credibility judgment owing to prejudice) from \emph{hermeneutical} injustice
(a structural gap in shared interpretive resources that leaves some
experiences unintelligible even to those who undergo them). The formal
apparatus of Section~\ref{p04:sec:formal} locates both.

\paragraph{Testimonial injustice as erasure.}
A partner asserts that sustained anticipatory care \emph{is} a contribution.
Under an operative interpretation $h$ that assigns such acts $\varepsilon$
(\S\ref{p04:def:interp}(c)), the assertion has no truth-maker in the record:
there is literally nothing in $\beta\circ h$ for the claim to be \emph{about}.
The credibility deficit then attaches not to the speaker's sincerity but to
the ontological standing of the asserted category. Empirically, this is the
familiar pattern in which unpaid affective and household labor---
disproportionately performed by women---is discounted in both private
negotiation and legal settlement \parencite{folbre2001}.

\paragraph{Hermeneutical injustice as kernel-blindness.}
More deeply: if the community's shared interpretive resources just \emph{are}
the operative $h$ (because the exchange grammar is the institutionally
supported one), then distinctions in $\ker h$ are not merely unrecorded but
\emph{unnameable}. The giver experiences a difference---between care and
appeasement, between a gift and a payment---for which the available
vocabulary has one word. This is hermeneutical injustice in the intimate
sphere, and it is mechanically coupled to alienation: the same
context-forgetting map that contracts practice (\S\ref{p04:sec:contraction}) also
impoverishes the concepts available for contesting the contraction. The
injustice is self-sealing.

% ------------------------------------------------------------
\section{From balance to recirculation: generative justice}
\label{p04:sec:generative}

\subsection{Eglash's generative justice}
Distributive justice asks how an already-produced surplus should be divided;
\textcite{eglash2016} asks instead whether value is allowed to \emph{circulate
back to, and remain within, the network of relations that generated it},
rather than being extracted from it. Eglash's examples---from commons-based
production to indigenous and DIY economies---share a topology: value flows in
\emph{cycles of recursive depth}, re-entering and compounding within the
generating loop, instead of exiting along an extraction path.\footnote{Paper
III of this series \parencite{huang2026vow} arrived at a cognate notion from the
normative side: generative justice as the maintenance of the conditions under
which each party's ongoing self-formation in the relation remains open and
uncaptured. The two senses converge in the present framework: the
recirculation of relational value (Eglash's sense) is the material condition
of ongoing co-formation (Paper III's sense), and the standing to contest the
terms is what keeps the recirculation from being unilaterally redirected.}

\subsection{The relation as the generative unit}
On a relational ontology, the unit of analysis in intimate giving is not the
individual party but the relation as a value-generating system: gifts are
moments in the relation's self-production, not transfers between pre-formed
endpoints. This re-founds the justice question. The ledger functional
$\beta$ evaluates \emph{node balance}---did each party's inflows match
outflows?---which presupposes exactly the per-symbol transfer reading whose
inadequacy Section~\ref{p04:sec:formal} established. The generative alternative
evaluates whether what the relation produces (material, affective,
\emph{and recognitional} value) recirculates.

\begin{definition}[Recirculation depth; sketch]
\label{p04:def:recirc}
Model a relation as a directed multigraph whose nodes include the parties
and the relation-level stock itself, with value-flow edges labeled by type
(material, affective, recognitional). For a flow history $F$, let $R(F)$ be
the proportion of generated value that lies on cycles returning to the
generating subgraph within the history, weighted by cycle depth
(number of re-entries). \emph{Generative justice} in the intimate sphere is
the requirement that $R$ be sustained across value types---in particular,
that recognitional flows not be systematically acyclic (recognition flowing
out of, but never back to, a contributor).
\end{definition}

\begin{remark}
$R$ deliberately treats \emph{recognition} as a first-class flow rather than
a residual. This is the formal repair of Section~\ref{p04:sec:epistemic}: an
act erased by $h$ ($\varepsilon$-mapped) cannot enter $\beta$, but it can
enter $R$, because $R$'s edge set is defined over the expressive practice,
not over its accounting image. A full treatment---axioms for $R$,
its relation to network-flow and to Eglash's recursive depth, and estimable
proxies---is deferred to a companion paper.
\end{remark}

\subsection{What the reframing dissolves}
The three harms now align as three faces of one error---applying $\beta$
where $R$ is the appropriate functional:
expressive alienation is the contraction of practice onto $\beta$'s domain
(\S\ref{p04:sec:contraction}); epistemic injustice is the unnameability of
$\ker h$ (\S\ref{p04:sec:epistemic}); distributive misdescription is the
evaluation of a cyclic generative system by an acyclic balance metric.

% ------------------------------------------------------------
\section{Implications: relational data and evidence-informed decisions}
\label{p04:sec:data}

The analysis bears directly on the utilization of intimate-relationship data
and on evidence-informed decision-making about relationships---an agenda this
author pursues in parallel empirical work.

\paragraph{Instrumentation imports a grammar.}
To measure is, by default, to implement an $h$: a relational data system that
counts gifts, chores, or messages as balanced inputs operationalizes the
ledger and thereby \emph{accelerates} the contraction of
\S\ref{p04:sec:contraction}---a relational instance of Goodhart's law. The
crowding-out literature (\S2.5) predicts, and Definition~\ref{p04:def:alienation}
explains, that such systems do not neutrally record the practice; they
re-form it.

\paragraph{Designing for the generative grammar.}
The constructive program is to encode $R$ rather than $\beta$:
(i) treat recognition as a logged, first-class flow, so that erased
contributions acquire truth-makers (repairing the testimonial deficit);
(ii) supply shared vocabulary for affective labor, expanding the
hermeneutical resources beyond the operative $h$;
(iii) evaluate the relationship by recirculation indicators (are
contributions of every type eventually re-entered, acknowledged,
reciprocated-in-kind-or-otherwise within the cycle?) rather than by balance;
(iv) preserve, in the data model itself, the context-dependence of meaning---
e.g.\ event records that link to history rather than free-standing valued
entries. Whether intimate-relationship data systems help or harm is therefore
not a question about data per se, but about \emph{which grammar the system
makes operative}.

% ------------------------------------------------------------
\section{Conclusion}
\label{p04:sec:conclusion}

Economics has not ignored the gift; it has interpreted it---as allocation,
implicit contract, signal, inefficient transfer, or priced incentive---and
each interpretation is a context-forgetting homomorphism into a ledger
language of strictly weaker expressive power. The semiotic and
psychoanalytic analysis identified what the interpretation loses---the
positional value of the signifier and the desire it attests---and the formal
analysis showed the loss to be architectural. Alienation is what happens when
the interpretation acquires normative force and practice contracts onto its
preimage; epistemic injustice is the unnameability of the distinctions the
map collapses; and distributive misdescription is the use of a balance
functional on a generative cycle. The repair is not to abandon formalization
but to formalize the \emph{right} object: the recirculation of generated
value---material, affective, recognitional---within the relation that
produces it.

Read alongside Paper III, the series' two diagnoses are one structure seen
from two sides. There, a law of love legislated unilaterally institutes
domination even when sincere; here, a language of love interpreted
unilaterally institutes alienation even when accurate. In both, the remedy is
not purer intention but shared authority over the terms---there, the
co-authorship of the vow; here, the keeping-alive of the expressive grammar
in which the gift can still say what it means. Future work will axiomatize
the recirculation functional, connect it rigorously to Eglash's recursive
depth, and test the design program of Section~\ref{p04:sec:data} against field
data from a deployed relational infrastructure.

% ------------------------------------------------------------
\section*{Acknowledgements}
\addcontentsline{toc}{section}{Acknowledgements}
This paper is the fourth in a series on the philosophy of intimacy and the
theory of justice. I thank Ron Eglash for correspondence on generative
justice and recursive depth, on which Section~\ref{p04:sec:generative} builds.

This series, and this paper above all, are dedicated to the one I love most
deeply---the ``girl of the forest and the mountains'' to whom the first paper of the series was already
devoted. We were born of the same hometown, yet it was in Japan that we met.
To have met her is the finest gift the universe has ever given me; and what I
give in return is the whole of my life itself, offered as a gift. It is from
her, and not from any theory, that the thesis of this paper was first
learned.

In the interest of transparency, I note that an AI assistant was used in
preparing this manuscript, as a tool for drafting, structuring, and refining
the argument and its prose; the ideas, commitments, and final judgements are
my own, and I take full responsibility for the content.

\begin{center}
\zh{愿天下有情人，赠不为偿，受不为债，情意长流而不竭。}
\end{center}

% ------------------------------------------------------------
